\documentclass[11pt, letterpaper]{article}
% ----- PAQUETES -----
\usepackage[utf8]{inputenc}
\usepackage[T1]{fontenc}
\usepackage[spanish]{babel}
\usepackage{graphicx}
\usepackage[margin=2.3cm]{geometry}
\usepackage{xcolor}
\usepackage{titlesec}
\usepackage{booktabs}
\usepackage{enumitem}
\usepackage{xurl}
\usepackage{hyperref}
\usepackage{fancyhdr}
\usepackage{amsmath}
\usepackage{amssymb}
\usepackage{array}
\usepackage{tikz}
\usetikzlibrary{arrows.meta, babel, calc}

% ----- COLORES -----
\definecolor{UnalGreen}{RGB}{0,153,51}
\definecolor{DarkGray}{RGB}{60,60,60}
\definecolor{StripA}{RGB}{200,60,60}
\definecolor{StripB}{RGB}{60,90,180}

% ----- SECCIONES -----
\titleformat{\section}
{\normalfont\large\bfseries\color{UnalGreen}}
{\thesection}{1em}{}[{\titlerule[0.8pt]}]
\titleformat{\subsection}
{\normalfont\normalsize\bfseries\color{DarkGray}}
{\thesubsection}{1em}{}
\titlespacing*{\section}{0pt}{8pt}{4pt}
\titlespacing*{\subsection}{0pt}{6pt}{2pt}

% ----- ENLACES -----
\hypersetup{
	colorlinks=true,
	linkcolor=UnalGreen,
	urlcolor=UnalGreen
}

% ----- ENCABEZADO -----
\pagestyle{fancy}
\fancyhf{}
\lhead{\footnotesize Mecánica y Ondas para Biociencias (1000023) -- Laboratorio}
\rhead{\footnotesize Semestre 2026-2}
\cfoot{\thepage}

% ----- LISTAS COMPACTAS -----
\setlist[itemize]{leftmargin=*, itemsep=2pt, topsep=2pt, parsep=0pt}
\setlist[enumerate]{leftmargin=*, itemsep=3pt, topsep=3pt, parsep=0pt}

% ----- CAJA DE NOTA -----
\newcommand{\notabox}[1]{%
	\begin{center}
		\fcolorbox{UnalGreen}{UnalGreen!8}{%
			\parbox{0.94\textwidth}{\small #1}%
		}
	\end{center}
}

\begin{document}
	
	% =========================================================
	% ENCABEZADO
	% =========================================================
	\begin{center}
		\vspace*{-0.8cm}
		\IfFileExists{imagenes/universidad-nacional-de-colombia-sede-bogota-seeklogo.png}{%
			\includegraphics[width=0.15\textwidth]
			{imagenes/universidad-nacional-de-colombia-sede-bogota-seeklogo.png}\\[0.1cm]
		}{}
		
		{\Large\bfseries Práctica 3: Primera Ley de Newton}\\[0.1cm]
		{\large\bfseries Equilibrio de fuerzas en una mesa de fuerzas}\\[0.1cm]
		{\normalsize Mecánica y Ondas para Biociencias (1000023)}
	\end{center}
	
	\vspace{0.2cm}
	\noindent
	\begin{tabular}{@{}ll@{}}
		Modalidad: & Grupal (3 a 4 integrantes) \\
		Entregable: & Bitácora 3 \\
		Fecha de sesión: & 06 de octubre de 2026 \\
	\end{tabular}
	
	\notabox{\textbf{Antes de venir al laboratorio:} lean la descripción del
		experimento y respondan las \emph{preguntas previas} de la
		Sección~\ref{sec:previas}. El quiz individual del
		inicio de la sesión sale de esas preguntas.}
	
	% =========================================================
	\section{Objetivos}
	% =========================================================
	\begin{itemize}
		\item Sumar fuerzas vectorialmente usando componentes rectangulares.
		\item Aplicar la primera ley de Newton para estudiar un sistema en
		equilibrio.
		\item Predecir teóricamente una masa de equilibrio y contrastarla con la
		medida experimentalmente.
	\end{itemize}
	
	% =========================================================
	\section{Materiales}
	% =========================================================
	\begin{itemize}
		\item Mesa de fuerzas con transportador graduado.
		\item Tres poleas ajustables.
		\item Juego de masas y portapesas.
		\item Regla.
	\end{itemize}
	
	% =========================================================
	\section{Descripción del experimento}
	% =========================================================
	
	Una mesa de fuerzas tiene, en su centro, un anillo del que salen tres
	cuerdas. Cada cuerda pasa por una polea ubicada en el borde de la mesa y
	cuelga de ella una masa $M_i$. El peso de esa masa, transmitido por la
	cuerda, actúa sobre el anillo como una fuerza $\vec{F}_i$ dirigida hacia la
	polea, en un ángulo $\theta_i$ que se lee directamente en la graduación de
	la mesa.
	
	\vspace{0.3cm}
\begin{center}
	\begin{tikzpicture}[scale=1.0,>=Latex]
		% mesa circular (más grande para dar espacio)
		\draw[thick] (0,0) circle (3.0);
		\foreach \ang in {0,30,...,330} {
			\draw[gray!60] ({3.0*cos(\ang)},{3.0*sin(\ang)}) --
			({3.18*cos(\ang)},{3.18*sin(\ang)});
		}
		
		% Etiquetas de los grados (bien alejadas)
		\node[font=\small, gray!60!black] at (3.6,0) {$0^\circ$};
		\node[font=\small, gray!60!black] at (0,3.6) {$90^\circ$};
		\node[font=\small, gray!60!black] at (-3.6,0) {$180^\circ$};
		\node[font=\small, gray!60!black] at (0,-3.6) {$270^\circ$};
		
		% eje de referencia punteado
		\draw[dotted, gray] (-3.0,0) -- (3.0,0);
		
		% anillo central
		\fill (0,0) circle (2.5pt);
		
		% --- FUERZAS (más cortas para no tocar los grados) ---
		% F1 (0°): termina en 2.5, etiqueta manual a la derecha
		\draw[-{Latex[length=3.5mm]}, very thick, StripA]
		(0,0) -- (2.5,0);
		\node[font=\large, StripA] at (2.0,0.38) {$\vec{F}_1$};
		
		% F2 (90°): termina en 2.5, etiqueta manual arriba
		\draw[-{Latex[length=3.5mm]}, very thick, StripB]
		(0,0) -- (0,2.5);
		\node[font=\large, StripB] at (0.35,1.9) {$\vec{F}_2$};
		
		% F3 (220°): termina en 2.5, etiqueta manual abajo-izquierda
		\draw[-{Latex[length=3.5mm]}, very thick, UnalGreen!70!black]
		(0,0) -- ({2.5*cos(220)},{2.5*sin(220)});
		\node[font=\large, UnalGreen!70!black] at (-1.0,-1.7) {$\vec{F}_3$};
		
		% --- ARCOS DE ÁNGULO (radios distintos, etiquetas libres) ---
		% Arco para theta_2 (radio 0.7)
		\draw[gray] (0.7,0) arc (0:90:0.7);
		\node[font=\large] at (0.35,0.35) {$\theta_2$};
		
		% Arco para theta_3 (radio 1.1, más grande para no encimarse)
		\draw[gray] (1.1,0) arc (0:220:1.1);
		\node[font=\large] at (-1.3,-0.35) {$\theta_3$};
		
	\end{tikzpicture}
\end{center}
	
	\noindent
	Si el anillo permanece quieto, las tres fuerzas se equilibran:
	$\vec{F}_1 + \vec{F}_2 + \vec{F}_3 = \vec{0}$.
	
	% =========================================================
	\section{Preguntas previas}\label{sec:previas}
	% =========================================================
	
	Respóndanlas \textbf{antes} de la sesión. No se califican, pero el quiz de
	entrada sale de aquí.
	
	\begin{enumerate}
		\item \emph{Predicción} (respondan antes de hacer cálculos): si
		$\vec{F}_1$ está en $\theta_1 = 0^\circ$ y $\vec{F}_2$ en
		$\theta_2 = 90^\circ$, con magnitudes parecidas, ¿en qué cuadrante
		debe apuntar aproximadamente $\vec{F}_3$ para que el anillo quede en
		equilibrio?
		
		\item Enuncien la primera ley de Newton. ¿Qué significa que un objeto
		esté en \emph{equilibrio}?
		
		\item Cada cuerda pasa por una polea sin fricción y cuelga de ella una
		masa $M_i$. Expliquen por qué la tensión de esa cuerda --y por tanto la
		fuerza $F_i$ que actúa sobre el anillo-- es igual a $M_i\,g$.
		
		\item Si tres fuerzas mantienen un punto en equilibrio,
		$\vec{F}_1+\vec{F}_2+\vec{F}_3=\vec{0}$, demuestren que esto exige que
		las sumas de las componentes en $x$ y en $y$ sean cero \textbf{por
			separado}.
	\end{enumerate}
	
	% =========================================================
	\section{Teoría}
	% =========================================================
	
	\subsection*{De la masa colgante a la fuerza sobre el anillo}
	
	Para una masa $M_i$ en reposo, colgando de una cuerda que pasa por una
	polea sin fricción, la tensión de la cuerda es igual al peso de la masa:
	
	\begin{center}
		\begin{tikzpicture}[scale=0.8]
			% borde de la mesa y polea
			\draw[thick] (-1.8,0) -- (1.3,0);
			\node[left, font=\scriptsize] at (-1.8,0.15) {hacia el anillo};
			\draw[thick] (1.5,0) circle (0.2);
			\draw[thick] (-1.5,0.05) -- (1.32,0.05);
			% hilo vertical y masa
			\draw[thick] (1.7,0) -- (1.7,-1.1);
			\draw[thick, fill=UnalGreen!20] (1.45,-1.1) rectangle (1.95,-1.5);
			\node[font=\scriptsize] at (1.7,-1.3) {$M_i$};
			\draw[-{Latex[length=2.5mm]}, thick] (1.7,-1.5) -- (1.7,-1.95)
			node[right, font=\scriptsize] {$M_i\,g$};
		\end{tikzpicture}
	\end{center}
	
	\noindent
	Así, la magnitud de cada fuerza sobre el anillo es $F_i = M_i\,g$, dirigida
	hacia la polea correspondiente, en el ángulo $\theta_i$ leído en la mesa.
	
	\subsection*{Condición de equilibrio en componentes}
	
	La condición vectorial $\vec{F}_1+\vec{F}_2+\vec{F}_3=\vec{0}$ equivale a
	que las componentes en $x$ y en $y$ se anulen por separado:
	
	\begin{equation}
		M_1g\cos\theta_1 + M_2g\cos\theta_2 + M_3g\cos\theta_3 = 0
		\label{eq:compx}
	\end{equation}
	\begin{equation}
		M_1g\sin\theta_1 + M_2g\sin\theta_2 + M_3g\sin\theta_3 = 0
		\label{eq:compy}
	\end{equation}
	
	\noindent
	La gravedad $g$ aparece en los tres términos de cada ecuación, así que se
	cancela. Despejando $M_3$ de la ecuación (\ref{eq:compx}):
	
	\begin{equation}
		M_{3,\text{teo}} = -\frac{M_1\cos\theta_1 + M_2\cos\theta_2}{\cos\theta_3}
		\qquad (\cos\theta_3 \neq 0)
		\label{eq:m3x}
	\end{equation}
	
	\noindent
	o, si $\theta_3$ está cerca de $90^\circ$ o $270^\circ$ (donde
	$\cos\theta_3\approx 0$), se despeja de la ecuación (\ref{eq:compy}) en su
	lugar:
	
	\begin{equation}
		M_{3,\text{teo}} = -\frac{M_1\sin\theta_1 + M_2\sin\theta_2}{\sin\theta_3}
		\qquad (\sin\theta_3 \neq 0)
		\label{eq:m3y}
	\end{equation}
	
	\notabox{Las ecuaciones (\ref{eq:m3x}) y (\ref{eq:m3y}) deberían dar el
		mismo resultado: son dos formas de leer la \textbf{misma} condición de
		equilibrio. Se usa la que evita dividir por un número cercano a cero.}
	
	% =========================================================
	\section{Procedimiento}
	% =========================================================
	
	\begin{enumerate}
		\item Monten la mesa de fuerzas con sus tres poleas.
		\item Coloquen $M_1 = 30$ g en $\theta_1 = 0^\circ$ y
		$M_2 = 30$ g en $\theta_2 = 90^\circ$.
		\item Ajusten la tercera polea con una masa $M_3$ y un ángulo
		$\theta_3$ hasta que el anillo quede \textbf{centrado y en reposo}.
		\item Registren $M_3$ y $\theta_3$, con sus incertidumbres.
		\item Repitan para completar la tabla.
	\end{enumerate}
	
	\notabox{El anillo está en equilibrio cuando queda exactamente centrado
		respecto a las tres poleas. Verifíquenlo desplazándolo levemente: debe
		regresar al mismo punto.}
	
	\vspace{0.2cm}
	\begin{center}
		\small
		$\Delta M = $ \rule{1.5cm}{0.4pt} g
		\qquad
		$\Delta\theta = $ \rule{1.5cm}{0.4pt} $^\circ$
		\quad (resolución de las masas y del transportador de la mesa)
		
		\vspace{0.3cm}
		\begin{tabular}{|c|c|c|c|c|c|c|}
			\hline
			Medida & $M_1$ (g) & $\theta_1$ & $M_2$ (g) & $\theta_2$
			& $M_3$ (g) & $\theta_3$ \\
			\hline
			1 & 30 & $0^\circ$  & 30 & $90^\circ$  & & \\ \hline
			2 & 30 & $0^\circ$  & 50 & $90^\circ$  & & \\ \hline
			3 & 50 & $45^\circ$ & 30 & $180^\circ$ & & \\ \hline
			4 & 30 & $45^\circ$ & 30 & $135^\circ$ & & \\ \hline
			5 & 30 & $60^\circ$ & 30 & $120^\circ$ & & \\ \hline
		\end{tabular}
	\end{center}
	
	% =========================================================
	\section{Análisis de resultados}
	% =========================================================
	
	\subsection{Masa teórica}
	
	Para cada medida, calculen $M_{3,\text{teo}}$ con la ecuación
	(\ref{eq:m3x}) o (\ref{eq:m3y}), según cuál evite dividir entre un coseno o
	un seno cercano a cero.
	
	\begin{center}
		\small
		\begin{tabular}{|c|c|c|c|c|}
			\hline
			Medida & $M_3$ medido (g) & $\theta_3$ medido
			& $M_{3,\text{teo}}$ (g) & \%Error \\
			\hline
			\rule{0pt}{3ex} 1 & & & & \\ \hline
			\rule{0pt}{3ex} 2 & & & & \\ \hline
			\rule{0pt}{3ex} 3 & & & & \\ \hline
			\rule{0pt}{3ex} 4 & & & & \\ \hline
			\rule{0pt}{3ex} 5 & & & & \\ \hline
		\end{tabular}
	\end{center}
	
	\noindent
	donde
	\[
	\%\text{Error} = \frac{|M_3 - M_{3,\text{teo}}|}{M_{3,\text{teo}}}
	\times 100\%
	\]
	
	\subsection{Comparación con la incertidumbre}
	
	Para cada medida, verifiquen si $M_{3,\text{teo}}$ cae dentro del intervalo
	$M_3 \pm \Delta M$ que ustedes midieron. Esta comparación es más honesta
	que el \%Error solo: un error del 5\,\% puede ser perfectamente aceptable
	si su incertidumbre también es de ese orden, o puede ser demasiado grande
	si su incertidumbre es mucho menor.
	
	% =========================================================
	\section{Conclusiones}
	% =========================================================
	
	Además de sus conclusiones propias, discutan:
	
	\begin{itemize}
		\item ¿En cuántas de las cinco medidas $M_{3,\text{teo}}$ cayó dentro
		de $M_3 \pm \Delta M$? ¿Qué dice esto sobre la calidad de sus ajustes?
		\item ¿Su predicción de la pregunta previa 1 resultó correcta?
		\item ¿Qué dificultades encontraron para dejar el anillo exactamente
		centrado? ¿Cómo afecta eso a la incertidumbre de $\theta_3$?
		\item ¿Qué aspectos del montaje (fricción en las poleas, lectura del
		ángulo, nivelación de la mesa) pueden explicar las diferencias entre
		teoría y experimento?
	\end{itemize}
	
	% =========================================================
	\section{Referencias}
	% =========================================================
	
	\begin{itemize}
		\item L. Agulles Perdós, \emph{Guías de Laboratorio: Mecánica y Ondas
			para Biociencias}, Departamento de Física, Universidad Nacional de
		Colombia, 2026. Laboratorio 6.
		\item C. Perilla, \emph{Guía de trabajo en el laboratorio: Todos al
			tiempo y equilibrados}, Universidad Nacional de Colombia, 2023.
		\item L. G. Chica G., \emph{Guías de Laboratorio de Física:
			Cinemática, Dinámica y Termodinámica}, Universidad Nacional de
		Colombia, 2003.
	\end{itemize}
	
	\vfill
	\noindent
	{\footnotesize \textit{Auxiliar docente: Miguel Andrés Perdomo Gutiérrez
			--- mperdomog@unal.edu.co}}
	
\end{document}